\begin{ouverture}[college4]{18}{Lumière : propagation et distances}{Partie 4 · Des signaux pour observer et communiquer}{Voir loin, est-ce voir le passé ?}
\objectif{Distinguer source primaire, diffusion et réception.}
\objectif{Exploiter le modèle du rayon et construire une ombre.}
\objectif{Relier distance, durée et vitesse de la lumière.}
\prerequis{Trajets, vitesse moyenne et organisation du Système solaire.}
\end{ouverture}
\pcPage{Le cours}{Lumière : propagation et distances}
\phantomsection\label{ess-18}
\essSection{1}{De la source à l’œil}
Une source primaire émet de la lumière ; un objet diffusant renvoie dans plusieurs directions une partie de la lumière reçue. Pour voir un objet, de la lumière issue de cet objet doit entrer dans l’œil. L’œil ne projette pas un rayon pour « toucher » l’objet. Un écran peut recevoir de la lumière sans la détecter comme un capteur électronique ; on précise le rôle de chaque système.\par
\pcFigure{figures/c18-ombre.pdf}{Pour une source ponctuelle, les rayons limites construisent une zone d’ombre derrière un objet opaque. Schéma géométrique simplifié.}
\essSection{2}{Un modèle de propagation}
Dans le vide et dans un milieu transparent homogène comme l’air calme à l’échelle du dispositif, la lumière se propage en ligne droite. Un rayon est une ligne orientée représentant sa direction et son sens. Il ne dessine pas la largeur réelle d’un faisceau. Un objet opaque peut intercepter certains trajets et créer une ombre. Les sources étendues peuvent aussi produire une pénombre.\par
\clearpage
\essSection{3}{Une vitesse immense mais finie}
Dans le vide, $c=299\,792\,458\unite{m/s}$ exactement ; pour nos calculs, $c\approx300\,000\unite{km/s}$. On utilise $d=c\Delta t$ ou $\Delta t=d/c$ en gardant les unités cohérentes. Dans un matériau transparent, la lumière peut se propager moins vite que dans le vide. Une distance astronomique peut donc correspondre à un délai observable, même si les délais du laboratoire sont minuscules.\par
\pcFigure{figures/c18-trajet.pdf}{Un trajet aller-retour mesure une distance totale égale à deux fois la séparation des objets.}
\essSection{4}{Année-lumière et observation}
Une année-lumière est la distance parcourue par la lumière dans le vide en un an : environ $9{,}46\times10^{12}\unite{km}$. Une étoile à quatre années-lumière est observée avec environ quatre ans de retard sur l’émission. Ce délai ne signifie pas qu’elle est forcément éteinte aujourd’hui. Les observations du Soleil ou d’une source intense exigent des moyens spécialisés ; aucun regard direct n’est proposé dans ce livre.\par
\begin{pcEncadre}{method}{Un exemple expliqué}Pour une distance Terre-Soleil arrondie à $150$ millions de kilomètres, on calcule : \[\Delta t=\dfrac{150\,000\,000}{300\,000}=500\unite s,\] soit $8\unite{min}\,20\unite s$. On observe donc le Soleil avec ce délai approximatif. Le calcul utilise une distance et une vitesse arrondies, sans prédire toutes les variations orbitales.\end{pcEncadre}
\pcPage{Expérimenter et retenir}{Lumière : propagation et distances}
\phantomsection\label{activite-18}
\begin{pcEncadre}{method}{Investigation en classe}\textbf{Question.} Le modèle rectiligne prédit-il l’ombre ? Avec une petite lampe scolaire, un objet opaque et un écran, tracer d’abord les rayons limites sur un schéma. Mesurer les distances relatives puis comparer la prévision à l’ombre observée. Déplacer l’écran et décrire l’effet. Une source étendue donne des bords moins nets : discuter la limite de l’approximation ponctuelle. \textbf{À rendre :} rayons orientés, zones éclairées et zone d’ombre. Ni laser dirigé vers les yeux ni observation du Soleil.\end{pcEncadre}
\pcSousTitre{Ce que je dois savoir expliquer}\begin{questions}
\item Un objet est vu lorsque sa lumière entre dans l’œil.
\item Un rayon est un modèle orienté de propagation.
\item L’ombre résulte de trajets interceptés par un objet opaque.
\item La vitesse de la lumière dans le vide est finie.
\item L’année-lumière est une unité de distance.
\end{questions}
\begin{pcEncadre}{attention}{Une erreur à éviter}Pour un signal réfléchi, la durée mesurée correspond souvent à l’aller et au retour : ne pas attribuer tout $c\Delta t$ à la seule distance source-obstacle.\end{pcEncadre}
\pcRepere{cor-activite-18}{Pistes pour l’investigation}\pcGrille
\clearpage\markright{Je démarre}
\rubrique{Je démarre}
\par\noindent\begin{minipage}{\linewidth}
\begin{exercice}[1]{Source ou diffuseur}
\begin{questions}
\item Une lampe allumée et une feuille éclairée ont-elles le même rôle lumineux ?
\item Lequel est une source primaire ?
\end{questions}
\end{exercice}
\end{minipage}
\par\noindent\begin{minipage}{\linewidth}
\begin{exercice}[1]{Le bon sens}
\pcFigureEx{figures/ex-vision.pdf}{Orienter le segment entre feuille et œil.}
\begin{questions}
\item Dans un dessin de vision d’une feuille, orienter le rayon entre feuille et œil.
\item L’œil émet-il ce rayon pour voir ?
\end{questions}
\end{exercice}
\end{minipage}
\par\noindent\begin{minipage}{\linewidth}
\begin{exercice}[1]{La trajectoire}
\begin{questions}
\item Comment se propage la lumière dans un milieu transparent homogène, selon le modèle ?
\item Un rayon possède-t-il un sens ?
\end{questions}
\end{exercice}
\end{minipage}
\par\noindent\begin{minipage}{\linewidth}
\begin{exercice}[1]{Une grande vitesse}
\begin{questions}
\item Quelle valeur arrondie de c utilise-t-on en km/s ?
\item Est-ce une vitesse infinie ?
\end{questions}
\end{exercice}
\end{minipage}
\par\noindent\begin{minipage}{\linewidth}
\begin{exercice}[1]{La zone d’ombre}
\begin{questions}
\item Qu’est-ce qui crée une ombre derrière un objet opaque ?
\item Un milieu transparent crée-t-il toujours la même ombre ?
\end{questions}
\end{exercice}
\end{minipage}
\clearpage\markright{Je m’entraîne}
\rubrique{Je m’entraîne}
\par\noindent\begin{minipage}{\linewidth}
\begin{exercice}[2]{Un trajet en une seconde}
\begin{questions}
\item Avec $c=300\,000\unite{km/s}$, quelle distance en $2\unite s$ ?
\item Et en $0{,}10\unite s$ ?
\end{questions}
\end{exercice}
\end{minipage}
\par\noindent\begin{minipage}{\linewidth}
\begin{exercice}[2]{Une durée solaire}
\begin{questions}
\item Calculer le délai sur $150$ millions de kilomètres avec la vitesse du cours.
\item Exprimer en minutes et secondes.
\end{questions}
\end{exercice}
\end{minipage}
\par\noindent\begin{minipage}{\linewidth}
\begin{exercice}[2]{Un aller simple lunaire}
\begin{questions}
\item Pour $384\,000\unite{km}$, calculer la durée aller de la lumière.
\item Calculer l’aller-retour idéal.
\end{questions}
\end{exercice}
\end{minipage}
\par\noindent\begin{minipage}{\linewidth}
\begin{exercice}[2]{Une étoile}
\begin{questions}
\item Une étoile est située à 12 années-lumière. Quel retard approximatif a sa lumière ?
\item Cela donne-t-il son état actuel ?
\end{questions}
\end{exercice}
\end{minipage}
\par\noindent\begin{minipage}{\linewidth}
\begin{exercice}[2]{Une unité à convertir}
\begin{questions}
\item On prend $1\unite{al}=9{,}46\times10^{12}\unite{km}$. Quelle distance pour 2 al ?
\item Quelle durée de propagation dans le vide ?
\end{questions}
\end{exercice}
\end{minipage}
\clearpage\markright{Je raisonne}
\rubrique{Je raisonne}
\par\noindent\begin{minipage}{\linewidth}
\begin{exercice}[3]{Une mesure réfléchie}
\pcFigureEx{figures/ex-echo-lumineux.pdf}{}
\begin{questions}
\item Un signal lumineux revient après $0{,}004\unite s$. Avec $c=300\,000\unite{km/s}$, calculer la distance de l’obstacle.
\item Pourquoi diviser par deux ?
\end{questions}
\end{exercice}
\end{minipage}
\par\noindent\begin{minipage}{\linewidth}
\begin{exercice}[3]{Une ombre qui grandit}
\pcFigureEx{figures/ex-ecrans.pdf}{Même source et même objet ; deux positions possibles de l’écran.}
\begin{questions}
\item Avec une petite source et un objet fixes, on éloigne l’écran derrière l’objet. L’ombre géométrique grandit-elle dans le modèle ?
\item Quelle construction justifie la réponse ?
\end{questions}
\end{exercice}
\end{minipage}
\par\noindent\begin{minipage}{\linewidth}
\begin{exercice}[3]{Le rayon invisible}
\begin{questions}
\item Dans l’air très propre, on voit une tache sur l’écran mais pas forcément le trajet du faisceau de côté. Pourquoi ?
\item Que change une légère diffusion ?
\end{questions}
\end{exercice}
\end{minipage}
\par\noindent\begin{minipage}{\linewidth}
\begin{exercice}[3]{Deux échelles}
\begin{questions}
\item Un laboratoire mesure une distance de $3\unite m$. Avec $c\approx3\times10^8\unite{m/s}$, calculer le délai aller.
\item Pourquoi un chronomètre manuel ne convient-il pas ?
\end{questions}
\end{exercice}
\end{minipage}
\par\noindent\begin{minipage}{\linewidth}
\begin{exercice}[3]{Une conclusion médiatique}
\begin{questions}
\item Un titre dit « Cette étoile est loin, donc elle est déjà morte ». Est-ce démontré ?
\item Quelle conclusion est justifiée par la distance seule ?
\end{questions}
\end{exercice}
\end{minipage}
\pcPage{Défi du chapitre}{Lumière : propagation et distances}
\phantomsection\label{defi-ex-18}
\pcDefiTitre{18}{Le miroir spatial}
\begin{questions}
\item Un dispositif reçoit un écho lumineux après $2{,}4\unite s$. Estimer la distance au réflecteur avec $c=300\,000\unite{km/s}$.
\item Un élève annonce 720 000 km. Expliquer son erreur.
\item Si le délai augmente, quelle conclusion est possible sous les mêmes hypothèses ?
\end{questions}
\pcGrille
